\documentclass[14pt,a4paper]{extarticle}

\usepackage[T2A]{fontenc}
\usepackage[utf8]{inputenc}
\usepackage[russian]{babel}
\usepackage{cmap}
\usepackage{geometry}
\usepackage{array}
\usepackage{longtable}
\usepackage{xcolor}
\usepackage{ragged2e}
\usepackage{listings}
\usepackage{enumitem}
\usepackage{hyperref}

\geometry{left=30mm,right=15mm,top=20mm,bottom=20mm}
\hypersetup{colorlinks=true,linkcolor=black,urlcolor=blue}
\setlength{\parindent}{1.25cm}
\setlength{\parskip}{0.2em}
\renewcommand{\arraystretch}{1.35}
\sloppy

\newcolumntype{L}[1]{>{\RaggedRight\arraybackslash}p{#1}}
\newcommand{\code}[1]{\texttt{\detokenize{#1}}}
\newcommand{\pathcode}[1]{{\footnotesize\nolinkurl{#1}}}

\lstdefinestyle{codestyle}{
    basicstyle=\ttfamily\small,
    breaklines=true,
    frame=single,
    columns=fullflexible,
    keepspaces=true,
    showstringspaces=false,
    backgroundcolor=\color{gray!5}
}

\begin{document}

\begin{titlepage}
    \begin{center}
        САНКТ-ПЕТЕРБУРГСКИЙ НАЦИОНАЛЬНЫЙ\\
        ИССЛЕДОВАТЕЛЬСКИЙ УНИВЕРСИТЕТ ИТМО

        \vspace{1.5cm}
        Дисциплина: Бэк-энд разработка

        \vspace{1.5cm}
        \textbf{Отчет}

        \vspace{0.5cm}
        Домашнее задание 6

        \vspace{0.5cm}
        \textbf{Настройка CI/CD для автоматического развёртывания backend-приложения}

        \vfill

        \begin{flushright}
            Выполнил:\\
            Цатинян Артём\\
            Группа БР1.2

            \vspace{0.8cm}
            Проверил:\\
            Добряков Д. И.
        \end{flushright}

        \vfill
        Санкт-Петербург\\
        2026 г.
    \end{center}
\end{titlepage}

\section{Задача}

В рамках домашнего задания требовалось настроить автоматическое развёртывание backend-приложения на удалённый сервер с использованием GitHub Actions или другой CI/CD-системы. Автодеплой должен запускаться при обновлении кода в выбранной ветке репозитория.

\section{Ход работы}

\subsection{Выбор подхода}

Для CI/CD была выбрана GitHub Actions, так как проект хранится в GitHub-репозитории. Развёртываемое приложение --- микросервисная версия сервиса бронирования ресторанов из ЛР2-ЛР4.

Во время ручного деплоя было выявлено, что сборка Gradle на сервере с 2 GB RAM работает слишком медленно и создаёт высокую нагрузку. Поэтому pipeline построен так, чтобы тяжёлая часть выполнялась на GitHub runner:

\begin{enumerate}[nosep]
    \item GitHub Actions получает код репозитория.
    \item На runner устанавливается Java 21.
    \item Выполняется сборка Spring Boot jar-файлов.
    \item Готовые jar-файлы архивируются.
    \item Архив передаётся на сервер по SSH/SCP.
    \item На сервере архив распаковывается, после чего systemd-сервисы перезапускаются.
\end{enumerate}

Такой подход экономит ресурсы VPS: сервер не компилирует Kotlin-проект, а только запускает готовые jar-файлы.

\subsection{Workflow GitHub Actions}

Файл workflow добавлен по пути:

\begin{lstlisting}[style=codestyle]
.github/workflows/deploy-lab2.yml
\end{lstlisting}

Workflow запускается при push в ветку \code{lab4-dz6}, а также вручную через \code{workflow_dispatch}.

\begin{lstlisting}[style=codestyle]
on:
  push:
    branches:
      - lab4-dz6
  workflow_dispatch:
\end{lstlisting}

\subsection{Сборка приложения}

В CI используется Java 21:

\begin{lstlisting}[style=codestyle]
- name: Set up Java
  uses: actions/setup-java@v4
  with:
    distribution: temurin
    java-version: "21"
\end{lstlisting}

Далее выполняется сборка:

В workflow рабочей директорией указан каталог \code{labs/lab2} внутри папки студента. Команда сборки:

\begin{lstlisting}[style=codestyle]
chmod +x ./gradlew
./gradlew --console=plain bootJar
\end{lstlisting}

После сборки создаётся архив с нужными артефактами:

\begin{lstlisting}[style=codestyle]
tar -czf lab2-jars.tar.gz \
  deploy-systemd.sh \
  identity-service/build/libs/identity-service.jar \
  catalog-service/build/libs/catalog-service.jar \
  booking-service/build/libs/booking-service.jar \
  review-service/build/libs/review-service.jar
\end{lstlisting}

\subsection{Передача артефактов на сервер}

Для загрузки архива используется action \code{appleboy/scp-action}. Он подключается к серверу по SSH и кладёт архив в \code{/tmp/lab2-jars.tar.gz}.

Для подключения используются secrets репозитория:

\begin{longtable}{|L{0.3\textwidth}|L{0.48\textwidth}|}
\hline
\textbf{Secret} & \textbf{Назначение} \\
\hline
\code{DEPLOY_HOST} & IP-адрес сервера: \code{186.246.13.204} \\
\hline
\code{DEPLOY_USER} & Пользователь SSH: \code{root} \\
\hline
\code{DEPLOY_PASSWORD} & Пароль SSH пользователя \\
\hline
\code{DEPLOY_PORT} & SSH-порт, необязательный secret, по умолчанию 22 \\
\hline
\code{DEPLOY_PATH} & Путь установки, необязательный secret, по умолчанию \pathcode{/opt/restaurant-booking-lab2} \\
\hline
\end{longtable}

\subsection{Перезапуск сервисов}

После загрузки артефактов запускается SSH-шаг:

\begin{lstlisting}[style=codestyle]
mkdir -p "$APP_DIR"
tar -xzf /tmp/lab2-jars.tar.gz -C "$APP_DIR"
cd "$APP_DIR"
sed -i 's/\r$//' deploy-systemd.sh
chmod +x deploy-systemd.sh
./deploy-systemd.sh
\end{lstlisting}

Скрипт \code{deploy-systemd.sh} создаёт и обновляет unit-файлы systemd для четырёх сервисов:

\begin{longtable}{|L{0.42\textwidth}|L{0.35\textwidth}|}
\hline
\textbf{Systemd unit} & \textbf{Сервис} \\
\hline
\pathcode{restaurant-booking-identity.service} & \code{identity-service} \\
\hline
\pathcode{restaurant-booking-catalog.service} & \code{catalog-service} \\
\hline
\pathcode{restaurant-booking-booking.service} & \code{booking-service} \\
\hline
\pathcode{restaurant-booking-review.service} & \code{review-service} \\
\hline
\end{longtable}

Сервисы включаются в автозагрузку и перезапускаются:

\begin{lstlisting}[style=codestyle]
systemctl daemon-reload
systemctl enable restaurant-booking-identity restaurant-booking-catalog restaurant-booking-booking restaurant-booking-review
systemctl restart restaurant-booking-identity
systemctl restart restaurant-booking-catalog
systemctl restart restaurant-booking-booking
systemctl restart restaurant-booking-review
\end{lstlisting}

\subsection{Проверка после деплоя}

В конце workflow выполняется smoke-check OpenAPI endpoints:

\begin{lstlisting}[style=codestyle]
curl -fsS http://localhost:8081/v3/api-docs >/dev/null
curl -fsS http://localhost:8082/v3/api-docs >/dev/null
curl -fsS http://localhost:8083/v3/api-docs >/dev/null
curl -fsS http://localhost:8084/v3/api-docs >/dev/null
\end{lstlisting}

Если хотя бы один сервис не запустился, команда \code{curl -fsS} завершится с ошибкой, и GitHub Actions пометит деплой как неуспешный.

\subsection{Результат}

После настройки сервер содержит рабочую ручную инсталляцию приложения, а pipeline автоматизирует повторный деплой при обновлении ветки \code{lab4-dz6}. Сервисы доступны по публичным адресам:

\begin{longtable}{|L{0.28\textwidth}|L{0.55\textwidth}|}
\hline
\textbf{Сервис} & \textbf{Swagger UI} \\
\hline
\code{identity-service} & \pathcode{https://ezhidze.ru/identity/swagger-ui/index.html} \\
\hline
\code{catalog-service} & \pathcode{https://ezhidze.ru/catalog/swagger-ui/index.html} \\
\hline
\code{booking-service} & \pathcode{https://ezhidze.ru/booking/swagger-ui/index.html} \\
\hline
\code{review-service} & \pathcode{https://ezhidze.ru/review/swagger-ui/index.html} \\
\hline
\end{longtable}

\section{Вывод}

В рамках домашнего задания был настроен CI/CD pipeline на GitHub Actions. Pipeline автоматически собирает микросервисное backend-приложение, передаёт готовые jar-файлы на сервер и перезапускает systemd-сервисы. Такой вариант учитывает ограниченные ресурсы VPS и переносит тяжёлую сборку на GitHub runner.

В результате обновление ветки \code{lab4-dz6} может автоматически приводить к обновлению работающего приложения на удалённом сервере.

\end{document}
